\chapter{Evaluation}\label{ch:analysis}

This chapter evaluates the \ac{GDI} pipeline against the three hypotheses
(H1, H2, H3) that operationalise the research questions from
\cref{ch:introduction}. It reports the design and results of each
experiment on the CPSCore case study, a complementary practitioner
session on a second industrial codebase and a discussion of threats to
validity, including several methodological risks caught during the
experiments themselves.

\section{Evaluation Design}\label{sec:eval-design}

\begin{itemize}
  \item \textbf{H1} (\cref{sec:h1}) compares reconstruction accuracy
    across conditions C1--C5 (\cref{sec:conditions}) on four scenarios
    (\cref{sec:scenarios}), all occurring naturally in CPSCore (S1, S2,
    S3, S4).
  \item \textbf{H2} (\cref{sec:h2}) compares guided against full
    instrumentation on structural compression, noise and coverage, with a
    blinded expert rating of the resulting diagrams.
  \item \textbf{H3} (\cref{sec:h3}) tests whether the reconstructed models
    improve an agent's performance on a scenario-comprehension task,
    relative to reasoning over raw source code.
\end{itemize}

\Cref{sec:rq-answers} draws the two research questions' answers together
from these experiments.
\Cref{sec:practitioner} reports a fourth, exploratory data point:
practitioner feedback from a second, much larger industrial codebase,
kept separate from H1--H3 since it was not run as a controlled or
pre-registered comparison.

\section{H1: Does Combining Static and Dynamic Evidence Improve Reconstruction Accuracy?}\label{sec:h1}

\textbf{Hypothesis.} Combining static analysis and dynamic instrumentation
via provenance-aware agreement (C5) improves reconstruction accuracy over
any single source (C1--C3) alone; raw set union (C4) does not, since it
can only inherit each source's false positives rather than cancel them.

\subsection{Metric}

For every condition and scenario, the generated \sysml sequence diagram
is compared against a manually authored reference diagram over
\emph{interaction edges} (unique ordered $\langle\text{sender},
\text{receiver}\rangle$ pairs):

\[
  P = \frac{|\text{gen} \cap \text{ref}|}{|\text{gen}|}, \quad
  R = \frac{|\text{gen} \cap \text{ref}|}{|\text{ref}|}, \quad
  F_1 = \frac{2PR}{P + R}
\]

Reference diagrams were authored before any condition was scored against
them, to avoid tuning the reference to a particular output.
This formulation follows the standard information-retrieval
convention~\cite{manning:ir}, mirrors its recent use for comparing
static against dynamically-sampled call graphs~\cite{helm:totalrecall},
and matches the convention used within the sequence-diagram
reconstruction literature itself~\cite{taniguchi:seqtrace}.

A second metric scores the same generated output against the same
reference for \emph{execution order}, which the edge-set metric above
discards entirely (two conditions can share an identical interaction
set while disagreeing completely on the order events fired). Let
$\text{gen}$ and $\text{ref}$ now denote \emph{ordered} event lists,
with repetitions preserved, and let $\mathrm{LCS}(\text{gen},
\text{ref})$ be their longest common subsequence -- the longest run of
events appearing in both lists in the same relative order, though not
necessarily contiguously:

\[
  P = \frac{|\mathrm{LCS}(\text{gen}, \text{ref})|}{|\text{gen}|}, \quad
  R = \frac{|\mathrm{LCS}(\text{gen}, \text{ref})|}{|\text{ref}|}, \quad
  F_1 = \frac{2PR}{P + R}
\]

Precision here penalises a condition for inserting extra events (e.g.\
uncorroborated static edges, or setup-phase trace events the
scenario-scoped ground truth excludes); recall penalises it for
dropping or misordering genuine ones. Both metrics are reported
throughout this section: the edge-set metric answers whether a
condition identifies the right \emph{interactions}; the LCS metric
answers whether it also gets their \emph{order} right.

\subsection{Results}\label{sec:h1-results}

\begin{table}[htb]
\centering
\caption{H1 per-scenario results (interaction edges): precision / recall
  / F1 by condition. All four scenarios occur naturally in CPSCore
  (\cref{sec:scenarios}).}
\label{tab:h1-per-scenario}
\small
\begin{tabular}{@{} l c c c c @{}}
\toprule
\textbf{Condition} & \textbf{S1} & \textbf{S2} & \textbf{S3} & \textbf{S4} \\
\midrule
C1: LLM-only        & 1.00/1.00/1.00 & 0/0/0.00       & .33/.33/.33 & 1.00/.33/.50 \\
C2: Static-only     & .50/1.00/.67   & 1.00/1.00/1.00 & 1.00/.33/.50 & .75/1.00/.86 \\
C3: Dynamic-only    & .67/1.00/.80   & 1.00/1.00/1.00 & 1.00/1.00/1.00 & .75/1.00/.86 \\
C4: Static+Dynamic  & .40/1.00/.57   & 1.00/1.00/1.00 & 1.00/1.00/1.00 & .60/1.00/.75 \\
C5: Intersection    & 1.00/1.00/1.00 & 1.00/1.00/1.00 & 1.00/.33/.50 & 1.00/1.00/1.00 \\
\bottomrule
\end{tabular}
\end{table}

\begin{table}[htb]
\centering
\caption{H1 per-scenario results (execution-order LCS metric, above): precision /
  recall / F1 by condition, over the same four scenarios.}
\label{tab:h1-per-scenario-seq}
\small
\begin{tabular}{@{} l c c c c @{}}
\toprule
\textbf{Condition} & \textbf{S1} & \textbf{S2} & \textbf{S3} & \textbf{S4} \\
\midrule
C1: LLM-only        & 1.00/1.00/1.00 & 0/0/0.00       & .58/.54/.56 & 1.00/.20/.33 \\
C2: Static-only     & .25/.13/.17    & 1.00/1.00/1.00 & 1.00/.08/.14 & .50/.40/.44 \\
C3: Dynamic-only    & .67/1.00/.80   & 1.00/1.00/1.00 & .77/1.00/.87 & .56/1.00/.72 \\
C4: Static+Dynamic  & .67/1.00/.80   & 1.00/1.00/1.00 & .77/1.00/.87 & .56/1.00/.72 \\
C5: Intersection    & 1.00/1.00/1.00 & 1.00/1.00/1.00 & 1.00/.15/.27 & 1.00/1.00/1.00 \\
\bottomrule
\end{tabular}
\end{table}

\begin{table}[htb]
\centering
\caption{H1 macro-averaged results across all four scenarios, interaction
  edges and execution-order LCS metric.}
\label{tab:h1-results}
\small
\begin{tabular}{@{} l ccc ccc @{}}
\toprule
& \multicolumn{3}{c}{\textbf{Interactions}}
& \multicolumn{3}{c}{\textbf{Sequence (LCS)}} \\
\cmidrule(lr){2-4} \cmidrule(lr){5-7}
\textbf{Condition} & P & R & F1 & P & R & F1 \\
\midrule
C1: LLM-only       & 0.583 & 0.416 & 0.458 & 0.646 & 0.434 & 0.473 \\
C2: Static-only    & 0.812 & 0.833 & 0.756 & 0.688 & 0.401 & 0.439 \\
C3: Dynamic-only   & 0.854 & \textbf{1.000} & \textbf{0.914} & 0.747 & \textbf{1.000} & \textbf{0.846} \\
C4: Static+Dynamic & 0.750 & \textbf{1.000} & 0.830 & 0.747 & \textbf{1.000} & \textbf{0.846} \\
C5: Intersection   & \textbf{1.000} & 0.833 & 0.875 & \textbf{1.000} & 0.788 & 0.817 \\
\bottomrule
\end{tabular}
\end{table}

F1 aggregation: macro-average across the four scenarios throughout.

\subsection{Discussion: Intersection Cancels Independent False Positives}\label{sec:h1-c4c2}

Three findings summarise the macro results in \cref{tab:h1-results}:
\textbf{C5 has the best precision of any condition, on both metrics}
(1.000 macro), never wrong when it commits to an edge; \textbf{C3 has
the best macro-F1 on interactions} (0.914, vs.\ C5's 0.875); and
\textbf{C4 ties C3 for the best macro-F1 on sequence} (0.846 each, vs.\
C5's 0.817). The rest of this section and \cref{sec:h1-c4c3,sec:h1-c5}
explain the mechanism behind each.

C4 (the raw union C2$\cup$C3) never beats the better of its two inputs
on interactions: its macro-F1 (0.830) sits strictly between C3's
(0.914) and C2's (0.756), because union can only \emph{inherit} each
source's false positives, never cancel them. C3 alone is the best
single-evidence condition, not merely a redundant subset of the static
graph: on S1, C2 carries two static over-approximation false positives
that the trace never exercises (P/R/F1~$=$~0.500/1.000/0.667); on S4 it
carries one (P/R/F1~$=$~0.750/1.000/0.857); on S3, C2's error is instead
a recall gap rather than a false positive (\cref{sec:h1-c4c3}); on S2
the two sources agree exactly. Recall is unaffected by S1/S4's false
positives -- these are precision losses only.

C5's perfect precision comes from responding to that same disagreement
in the opposite direction on S1, S2 and S4: since the two sources agree
on every true edge on each and disagree only on their own distinct
false ones -- C2's static over-approximation, C3's unfiltered
test-fixture-setup call in the trace -- keeping only confirmed edges
removes both sets of false positives at once:

\begin{itemize}
  \item S1 -- C2 additionally reports \texttt{CPSLogger::instance()}
    and one further static-only edge; C3 additionally reports the
    \texttt{Aggregator.add} test-setup call. C4 inherits both
    (P/R/F1~$=$~0.400/1.000/0.571); C5 keeps only what both confirm
    (P/R/F1~$=$~1.000/1.000/\textbf{1.000}).
  \item S2 -- C2 and C3 already agree exactly on the single genuine
    edge, so C4 and C5 coincide, both reaching
    P/R/F1~$=$~1.000/1.000/\textbf{1.000}.
  \item S4 -- C2 additionally reports the same static-only edge as on
    S1; C3 again reports the \texttt{Aggregator.add} setup call. C4
    inherits both (P/R/F1~$=$~0.600/1.000/0.750); C5 again reaches
    P/R/F1~$=$~1.000/1.000/\textbf{1.000}.
\end{itemize}

Static analysis has no way to scope itself to a single scenario, which
is the over-approximation problem that motivates requiring agreement
rather than simply pooling both sources' output; the trace has the
complementary problem of carrying setup-phase noise forward unfiltered.
Requiring agreement cancels both at once wherever each source's own
error is independent of the other's -- the case on three of the four
scenarios in this suite, all found-in-the-wild rather than constructed
for this evaluation (\cref{sec:scenarios}). C5's precision stays perfect
on the fourth (S3) too, but its recall does not -- which is why C3, not
C5, ends up with the best interactions F1 overall (\cref{sec:h1-c4c3}).

\subsection{Discussion: S3 and the Limits of Intersection}\label{sec:h1-c4c3}

S3 is the one scenario where C5 does not dominate: it collapses to
P/R/F1~$=$~1.000/0.333/0.500 (\cref{tab:h1-per-scenario}) -- precision
still perfect, but recall gone. The cause is a genuine recall gap, not a
disagreement intersection can resolve: one of S3's three genuine edges
(\texttt{Utilities}$\to$\texttt{Logging:stream}) is dispatched through a
\texttt{std::function}-wrapped thread with no resolvable static call
site (\cref{sec:scenarios}), so C2 recovers only one of three edges
while C3 recovers all three from the trace. Since intersection can
never exceed the smaller of its two inputs, it inherits C2's false
negative rather than C3's true positive. Aggregated across all four
scenarios (\cref{tab:h1-results}), this single miss is enough to pull
C5's interactions F1 (0.875) below C3's (0.914), even though C5's
precision (1.000 macro) never slips.

This is the general limit of agreement-based combining: it is a
precision filter, not a recall booster. Where the two sources'
\emph{errors} are independent -- S1, S2, S4 -- requiring agreement
strictly dominates both single sources at once. Where one source has a
genuine, structural blind spot the other does not share -- S3's
indirect dispatch -- intersection is bounded by the weaker source's
recall instead.

\subsection{Discussion: Where C3 and C4 Edge Out C5 on F1}\label{sec:h1-c5}

On interactions, S3's recall gap costs C5 far more than its occasional
false positives cost C3: C5 is perfect on three of the four scenarios
(\cref{tab:h1-per-scenario}) but collapses to 0.500 on S3; C3 is never
perfect on precision (0.667--0.750 on S1 and S4) but is perfect on S3,
where it observes all three genuine edges directly.

On sequence, C4 joins C3 at the top (0.846 each, vs.\ C5's 0.817) --
not by coincidence: C4's sequence artefact is a literal copy of C3's
trace ordering, since C2 contributes no genuine execution-order
information, only a depth-based heuristic; every
per-scenario sequence score in \cref{tab:h1-per-scenario-seq} is
identical between the two conditions. S3's recall gap costs C5 even
more here, because it removes not just one edge but nearly the entire
ordered sequence (seq~F1~$=$~0.267 on S3 alone,
\cref{tab:h1-per-scenario-seq}): dropping an edge from a \emph{set}
costs one false negative, but dropping it from a \emph{sequence} can
remove most of the ordered list at once if that edge recurs.

Which condition to prefer in practice therefore depends on whether
false positives or false negatives matter more downstream: C5
guarantees every reported edge is real; C3 (and, on sequence, C4)
guarantee no genuine edge on a well-covered scenario is missed, at the
cost of occasionally reporting one that is not scenario-relevant.

C2 is markedly worse on the sequence metric than on interactions (0.439
vs.\ 0.756 macro-F1): its depth-ordered heuristic (shorter call-graph
path approximates earlier execution) is a weak
proxy for actual execution order -- a structural limitation, not
scenario-specific noise.

\section{H2: Does Guided Instrumentation Improve Scenario Focus?}\label{sec:h2}

\textbf{Hypothesis.} Code-graph-guided instrumentation
(\cref{sec:guided-vs-full}) suppresses irrelevant execution detail
relative to full, unguided instrumentation.

H2 has a quantitative arm (structural compression, noise, coverage, and
a complementary reconstruction-accuracy check reusing the H1 conditions
and metrics) and a qualitative arm (blinded expert rating). All were
needed: H2 is fundamentally a claim about human-perceived usability,
which the quantitative metrics alone cannot establish.

\subsection{Metrics}

Let $G_{\text{full}}$, $G_{\text{guided}}$ denote the property-graph
subsets from full and guided instrumentation, and $R$ the
scenario-relevant reference node set (shared with H1):

\[
  \text{NodeReduction} = 1 - \frac{|V(G_{\text{guided}})|}{|V(G_{\text{full}})|},
  \qquad
  \text{EdgeReduction} = 1 - \frac{|E(G_{\text{guided}})|}{|E(G_{\text{full}})|}
\]

\[
  \text{Noise}(G) = \frac{|V(G) \setminus R|}{|V(G)|},
  \qquad
  \text{Coverage}(G) = \frac{|V(G) \cap R|}{|R|}
\]

Noise and Coverage are the retained-set analogues of $1-\text{Precision}$
and Recall~\cite{manning:ir}; the noise concept is adapted from its use
in architecture recovery evaluation~\cite{constantinou:noise}, applied
here to instrumentation scoping rather than clustering. A configuration
that reduces noise by discarding relevant nodes along with irrelevant
ones is not an improvement; coverage rules that failure mode out.

\subsection{Quantitative Results}\label{sec:h2-results}

\begin{table}[htb]
\centering
\caption{H2 aggregate results: guided vs.\ full instrumentation (mean across
  S1--S4).}
\label{tab:h2-results}
\begin{tabular}{@{} l c c c @{}}
\toprule
\textbf{Metric} & \textbf{Full} & \textbf{Guided} & \textbf{Reduction} \\
\midrule
Nodes (lifelines)     & 6.0 & 3.0 & $-50\%$ \\
Interaction edges     & 9.0 & 2.7 & $-70\%$ \\
Coverage              & 1.000 & 1.000 & $0\%$ \\
Noise                 & 0.704 & 0.000 & $-100\%$ \\
\bottomrule
\end{tabular}
\end{table}

Guided instrumentation reduced retained nodes by 50\% and interaction
edges by 70\%, driving noise to 0 while holding coverage at 1.000: every
scenario-relevant interaction survived the reduction. Full-trace false
positives are real runtime events from non-primary components active in
the same test run (\eg \texttt{Aggregation$\to$Logging:stream}),
architecturally valid but irrelevant to the scenario at hand.

The quantitative arm above is therefore scoped
to S1--S4, where full instrumentation actually produces excess events
for guided instrumentation to remove.

\subsection{Qualitative Arm: Blinded Expert Rating}\label{sec:h2-qual}

Structural compression alone cannot establish that guided instrumentation
produces a \emph{more usable} diagram. That is a claim about human
perception. A blinded A/B protocol tested this directly: six external
domain engineers, with no prior knowledge of the codebase, independently
rated diagrams from both configurations without being told which was
which, using a rating form and a private decode key (held separately) to
keep the assignment blind, with the A/B label randomised per scenario.
All six returned completed forms scoring S1--S3 on accuracy, readability
and practical usefulness, plus free-text comments and an overall
preference.

\paragraph{Results}

\begin{table}[htb]
\centering
\caption{H2 qualitative results: mean blinded expert ratings (1--5 scale,
  $n=6$ raters) for full vs.\ guided instrumentation, per scenario and
  rating dimension.}
\label{tab:h2-qual-results}
\begin{tabular}{@{} l cc cc cc @{}}
\toprule
 & \multicolumn{2}{c}{\textbf{Accuracy}} & \multicolumn{2}{c}{\textbf{Readability}} & \multicolumn{2}{c}{\textbf{Usefulness}} \\
\textbf{Scenario} & Full & Guided & Full & Guided & Full & Guided \\
\midrule
S1 & 2.33 & 4.67 & 2.33 & 4.83 & 3.17 & 4.50 \\
S2 & 2.33 & 4.17 & 2.33 & 4.83 & 2.67 & 2.83 \\
S3 & 2.50 & 5.00 & 2.67 & 5.00 & 3.00 & 4.67 \\
\midrule
\textbf{Mean} & \textbf{2.39} & \textbf{4.61} & \textbf{2.44} & \textbf{4.89} & \textbf{2.94} & \textbf{4.00} \\
\bottomrule
\end{tabular}
\end{table}

Guided instrumentation was never rated below full instrumentation on
accuracy or readability by any rater on any scenario (34 of 36
comparisons strictly preferred, 2 exact ties). Across all
$6 \times 9 = 54$ rater--scenario--dimension comparisons, guided was
rated strictly higher in 47, tied in 6, and lower in exactly 1. Averaged
across all scenarios and dimensions, guided scored 4.50 against 2.59 for
full instrumentation.

The larger panel did not change the direction of the effect but
surfaced heterogeneity a two-rater sample could not. Usefulness is
unanimous for guided on S1 and S3 but genuinely split on S2, the one
scenario where guided instrumentation collapses seven repeated
\texttt{Configuration}$\to$\texttt{Logging} property-read calls into a
single edge: three raters tied it, one preferred full, two preferred
guided. Free-text comments suggest this reflects a real trade-off:
collapsing repeated reads discards a \emph{count} some raters consider
useful for debugging, independent of accuracy or readability. Full
per-rater decomposition and comments are in
\cref{sec:appendix-h2-perrater}.

Taken together, guided instrumentation compresses the graph without
losing scenario-relevant nodes (quantitative arm) and is rated more
accurate and more readable by every one of six domain experts, on every
scenario (qualitative arm). The one qualification is usefulness on S2,
where the reduction discards a property-read count several raters found
independently useful for debugging.

\section{H3: Do Reconstructed Models Improve Task Performance?}\label{sec:h3}

\textbf{Hypothesis.} Agents equipped with reconstructed structural and
sequence diagrams perform scenario-comprehension tasks more accurately
than agents relying only on source code and traces.

H3 was scoped to a single downstream task (Scenario Comprehension) to
keep the evaluation tractable. A single holistic experiment isolated
source code, execution traces, and a rendered diagram individually and
in combination, across all four scenarios, S1--S4.

\subsection{Protocol}

Ground truth for this rescoring is the same automatically derived,
scenario-scoped interaction set used throughout H1
(\texttt{gt\_interactions.json} per scenario, \cref{sec:scenarios}),
converted into expected components and expected interactions. Rescoring against the same ground
truth H1 uses keeps the two hypotheses consistent with each other. For
each scenario, an agent was prompted to identify components, enumerate
interactions and explain the execution flow under one of five
conditions:

\begin{description}
  \item[Source only.] Source snippets only.
  \item[Trace only.] Runtime trace subset only.
  \item[Diagram only.] A rendered \sysml v2 diagram (structural and
    sequence, \texttt{gen\_sysml.py}, \cref{sec:sysml-generation}) only.
  \item[Source + Trace.] Both.
  \item[Source + Trace + Diagram.] All three.
\end{description}

Component and interaction recall apply the same formulation as H1: a
response is credited for a ground-truth interaction only if the source
component, target component and bare method name are all genuinely
attributed to that specific interaction, not merely present somewhere
in the response text. Completeness is scored 1--5 against
\texttt{scoring/rubric.md} by a single non-blind rater applying the
ground truth above.

\subsection{Holistic Comparison}
\label{sec:h3-holistic}

\begin{table}[htb]
\centering
\caption{H3 holistic comparison: five conditions across all four scenarios, S1--S4.}
\label{tab:h3-holistic}
\small
\begin{tabular}{@{} l c c c @{}}
\toprule
\textbf{Condition} & \textbf{Component recall} & \textbf{Interaction recall} & \textbf{Completeness (1--5)} \\
\midrule
Source only               & 0.833 & 0.444 & 3.00 \\
Trace only                 & \textbf{1.000} & 0.778 & 4.00 \\
Diagram only               & \textbf{1.000} & 0.778 & 4.25 \\
Source + Trace              & \textbf{1.000} & \textbf{1.000} & 4.75 \\
Source + Trace + Diagram    & \textbf{1.000} & \textbf{1.000} & \textbf{5.00} \\
\bottomrule
\end{tabular}
\end{table}

Diagram-inclusive conditions are never worse than their non-diagram
counterparts, and are often better. Diagram only clearly beats source
only (0.778 vs.\ 0.444 interaction recall) and ties trace only, while
adding the diagram to source+trace lifts completeness further (5.00
vs.\ 4.75) at no cost to recall. Source only is the weakest condition
throughout, missing a component and interactions on S1 and S4 that
trace or diagram evidence make explicit.

The one shared limitation is S3: source only, trace only and diagram
only each misattribute one of two \texttt{stream} interactions to
\texttt{Synchronization} rather than their true source component,
capping single-evidence interaction recall there at 1 of 3. Only
combining source and trace resolves it, by letting the agent
cross-reference the trace's specific caller against the source code's
specific class -- which is why source+trace and source+trace+diagram
are the only conditions to reach perfect recall.

\subsection{H3 Conclusion}

H3 is supported: diagram-inclusive conditions are never worse than
their non-diagram counterparts and clearly outperform source only, the
weakest condition on every measure. Diagram only ties trace only on
recall but edges ahead on completeness, and adding the diagram to
source+trace raises completeness further without cost to recall. The
one qualification is S3, where a diagram alone -- like source or trace
alone -- cannot resolve which component actually originates a shared
interaction name; only combining source and trace does that.
\cref{sec:rq-answers} revisits \rqthree in light of this result.

\section{Practitioner Validation: A Second Industrial Codebase}\label{sec:practitioner}

H1--H3 validate the approach in depth on one system (CPSCore) under
controlled conditions. This section reports a complementary, exploratory
data point. \emph{Renaissance MCP}, the tool built in this project,
exposes the Renaissance static-extraction backend and graph-querying
approach as an MCP integration inside GitHub Copilot, and was trialled
by practitioners at an industrial partner on a second \Cpp\ codebase
several orders of magnitude larger than CPSCore.

This is not a repeat of H1/H2/H3: it is an unblinded, semi-structured
hands-on session and interview study with four engineers, exercising
only the static-extraction and graph-querying capability. No runtime
instrumentation or dynamic analysis was involved, so it speaks to the
static half of the approach only, and is reported because it bears
directly on external validity: whether the static-analysis capability
holds up on a second, much larger system. Codebase and defect details
are kept generic to avoid disclosing partner-internal specifics.

\subsection{Session Design}

Each session had four parts: tool installation, an \emph{estimation
game} (guessing then checking quantitative facts about the codebase), a
\emph{scenario exercise} investigating a real, previously-resolved
engineering issue, and a wrap-up discussion on the tool's diagram
outputs.

\begin{table}[htb]
\centering
\caption{Estimation-game answers, illustrating codebase scale on the
  second industrial system.}
\label{tab:renaissance-scale}
\begin{tabular}{@{} l r @{}}
\toprule
\textbf{Query} & \textbf{Answer} \\
\midrule
Method declarations on the largest class (a test helper) & 923 \\
Header files not included elsewhere (excl.\ externals)   & 9,606 \\
Inclusions of the most-included header file              & 2,053 \\
Calls to the most-called function                         & 23,846 \\
Events declared by a representative event-heavy interface & 118 \\
Call sites of a representative header-declared method      & 0 (unused) \\
\bottomrule
\end{tabular}
\end{table}

The scale here (\cref{tab:renaissance-scale}) contrasts with CPSCore's
2,574 static interactions and 121 components (\cref{tab:impl-summary}).

\subsection{Scenario Exercise: A Real Engineering Issue}\label{sec:scenario-exercise}

The exercise centred on a previously-resolved defect chosen for matching
this thesis's target problem class: an implicit, event-driven
interaction invisible without reconstructing the call/event chain.

\textbf{Defect summary (generalised, anonymised).} A control service
exposes interfaces to a client via proxy objects created on connect. On
disconnect, most proxies are freed; one was not, due to a missing reset
in connection-handling code. After a later restart, the stale proxy
silently stopped delivering change-notification events. The defect was
visible only by tracing the event path from the reconnect handler,
through the stale proxy, to the downstream handler.

\textbf{Task.} Participants used the tool to answer three questions
about the already-fixed codebase: which files changed, which test
exercises the fix, and where the proxy objects are used. All three are
answerable from the static call/containment graph alone, making this a
check on the same static capability underlying \rqone\ and \rqtwo, now
on an unfamiliar, industrial-scale codebase. Participants correctly
identified the modified file, the relevant test, and the proxy's
construction/release sites, the components underlying the actual fix.

\subsection{Interview Findings}\label{sec:interview-findings}

Four practitioners were interviewed: a software designer (P1), a product
owner for an embedded control service (P2), a senior architect for
inter-component synchronisation (P3), and an algorithm engineer (P4).
Interviews covered general usefulness and six diagram types: class
structure (class diagram, polymetric view), object interaction (call
graph, dependency diagram), and system overview (structural, sequence
diagram).

All four use AI coding assistants daily and saw value in
architecture/dependency information. P3 reported completing an
investigation in a five-minute session that had previously taken roughly
five months, a striking self-report, but a single, unverified,
retrospective comparison, reported as anecdote rather than measured
effect. P2 noted the graph-based tool located a stub implementation
whose concrete definition lived outside the visible workspace, a failure
mode of code-search-based assistants. P4 found it most useful for
onboarding onto unfamiliar code and locating coverage gaps.

Feedback on diagrams converged on a pattern. Dependency and structural
diagrams were consistently rated useful, and in P3's case judged cleaner
than the team's existing modelling-tool output. Sequence diagrams were
valued for surfacing unexpected race conditions and cross-component
effects (P1, P3, P4), but all three found event/message labels (\eg a
generic \texttt{Subscribe}) too abstract without inline detail, and
requested source class/line number on hover. The class diagram and
polymetric view drew the most mixed feedback: three of four preferred
the plain class diagram, and the polymetric view's size-based encoding
was repeatedly flagged as measuring the wrong thing (P2, P4 suggested
cyclomatic complexity or requirements traceability instead). Call graphs
were the most polarising output: redundant with existing build-system
views for P4, not what P3 expected, while P2 rated the equivalent
dependency diagram positively for indicating ``hops'' between
components.

Three requests recurred across interviews, independent of diagram type:
interactivity (static diagrams compared unfavourably to tools allowing
search/filter/collapse, P1--P3); non-code context such as timing,
deployment configuration and product variant (P2, P4), a direct
instance of the construct-validity limitation already noted for this
thesis's own evaluation, since the graph
only contains what was extracted; and invocation friction, with some
participants needing unusually explicit phrasing to trigger the tool
(P1, P4).

\subsection{Relationship to This Thesis's Contribution}

The class diagram and polymetric view are contributed by other \ac{GDI}
members and fall outside this thesis's scope; they are reported only as
context. The findings most relevant here are the positive reception of
the dependency and structural diagrams (descended from the interface
dependency table, \cref{sec:idt}) and the actionable critique of
sequence-diagram labelling (relevant to \sysml generation,
\cref{sec:sysml-generation}). The hover-detail request is a low-effort
extension of the existing \texttt{sourceLineNumber} property already
carried on \texttt{CppCalls} edges (\cref{sec:utility}), noted as future
work (\cref{sec:future-work}).

The proxy-leak defect (\cref{sec:scenario-exercise}) is an interaction
architecturally present in source but only exercised under a specific
failure condition (reconnect after a missed release). Participants
resolved it using only the static graph, the same evidence source as
C2, which recovers architecturally-present interactions regardless of
whether they were exercised in any particular run
(\cref{sec:h1-c4c2}). This is informal, but convergent, evidence that
static extraction alone remains useful for fault-path investigation on a
second, much larger codebase.

The same architecture also motivates, without measuring, a prediction
about combining on this system (\cref{sec:cpscore-scale,sec:h1-c4c2}).
S1 and S4 show combining helps not because static and dynamic analysis
each recover a different genuine edge the other misses, but because
each source's own errors are independent, so requiring agreement (C5)
cancels them without discarding anything both sources confirm
(\cref{sec:h1-c4c2}); raw union (C4), by contrast, only ever inherits
more error. Unlike CPSCore, this codebase uses proxy objects, virtual
interfaces (118 events on one representative interface) and dynamic
connect/disconnect, conditions under which static analysis is more
likely to over- or under-approximate on its own -- exactly the setting
where a second, independently-erring source has more disagreement for
intersection to cancel. P2's stub-resolution observation, where the
graph-based tool located a definition invisible to code search, is
consistent with static analysis carrying its own blind spots here that a
disagreeing dynamic source could catch; but the session collected no
dynamic evidence, so it cannot confirm intersection would outperform C2
alone on this codebase, only that C2 alone already found something a
different tool missed. That mechanism was measured only on CPSCore, via
S1 and S4 (\cref{sec:h1-c4c2}). The industrial architecture motivates
the prediction that intersection would help here too; testing the
false-positive-cancellation mechanism on a system of this scale and
dispatch complexity is left to future work (\cref{sec:future-work}).

\section{Discussion: Cross-Cutting Pattern}\label{sec:discussion}

A single pattern runs through the H1 results: additional evidence helps
when it disagrees \emph{productively} with the primary source, by
independently confirming which of each source's edges are genuine, and
is otherwise redundant or actively harmful.

\begin{description}
  \item[H1.] raw union (C4) never beats the better single source,
    since it only inherits false positives. Agreement-based combining
    (C5) beats both where the sources' \emph{errors} are independent
    (S1, S2, S4), but is bounded by the weaker source's recall where
    one has a genuine gap the other doesn't share (S3,
    \cref{sec:h1-c4c3}).
  \item[H2.] Guided instrumentation removes noise without losing
    coverage, and the panel confirms this translates into perceived
    usability -- except on S2, where the removed information was itself
    useful for debugging.
  \item[H3.] Diagram-inclusive conditions beat source alone and are
    never worse than their non-diagram counterparts, but the diagram
    alone still shares source and trace's S3 attribution error; only
    combining source and trace resolves it, by letting the agent
    cross-reference the trace's specific caller against the source
    code's specific class (\cref{sec:h3-holistic}).
  \item[Practitioner session.] The proxy-leak defect was resolved using
    only the static graph; with no dynamic evidence collected, this
    speaks to static extraction's standalone value, not to the
    intersection mechanism itself (\cref{sec:practitioner}).
\end{description}

The pipeline's precision benefit is therefore conditional on sources
failing independently, not uniform -- useful for engineers deciding
whether and how to combine two extraction methods. H3 adds a downstream
caveat: a reconstructed diagram summarises what the pipeline concluded,
but attributing an interaction to its specific source component still
calls for cross-referencing source and trace directly.

This generates a testable prediction: on a codebase with cross-module
virtual dispatch (unlike CPSCore), static analysis would carry
independent errors on more scenarios, giving intersection more to
cancel. The practitioner session's codebase (proxy-based, dynamic
service binding) looks like such a system, and the engineers' difficulty
resolving stub implementations behind indirection is suggestive, if not
conclusive, evidence of this.

\section{Answering the Research Questions}\label{sec:rq-answers}

\subsection*{\rqone: Hybrid vs. Individual Sources}

\textit{Does combining a static call graph with runtime execution traces
reconstruct scenario-specific Interaction Models more accurately than
either source alone?}

Combining improves reconstruction \emph{conditionally}, not universally:
only agreement (C5~$=$~C2$\cap$C3), not raw union (C4), helps. On S1,
S2 and S4, intersection reaches perfect P/R/F1
(1.000/1.000/1.000), against 0.500--0.857 for single sources and
0.571--0.750 for union (\cref{sec:h1-c4c2}); on S3 it is bounded by
C2's recall gap on a \texttt{std::function}-dispatched call C3 observes
directly (\cref{sec:h1-c4c3}). S1 and S4 -- found-in-the-wild instances
of the same test, not constructed for this evaluation -- show the
mechanism holds without special construction, and union is never
protective on any scenario. Whether this generalises beyond CPSCore
depends on whether a codebase's static and dynamic errors are
independent, which the practitioner session (\cref{sec:practitioner})
can motivate but not confirm.

\subsection*{\rqtwo: Guided Instrumentation}

\textit{Does guiding trace collection with the static call graph remove
reconstruction noise without losing scenario coverage?}

Applying the scenario scope \emph{before} instrumentation reduces
retained nodes by 50\% and edges by 70\% while holding coverage at
1.000. The blinded expert rating confirms this translates into
perceived usability, with one qualification: collapsing repeated reads
on S2 discards a count several raters found useful for debugging
(\cref{sec:h2-qual}).

\subsection*{\rqthree: Sufficiency for Comprehension}

\textit{How well does a reconstructed \sysml Interaction Model
consolidate static and dynamic evidence into a single artefact sufficient
for system-behaviour comprehension?}

H3 (\cref{sec:h3}) supports the diagram's value but not its strict
sufficiency: diagram-inclusive conditions beat source only and are
never worse than their non-diagram counterparts, yet the diagram alone
still ties trace only and shares its S3 attribution error
(\cref{sec:h3-holistic}). Combining source and trace resolves that
case; adding the diagram on top raises completeness further at no cost
to recall. A reconstructed diagram is therefore a useful complement to
the underlying evidence, not a sufficient replacement for it.

\section{Engineering Utility}\label{sec:utility}

Beyond precision/recall, a tool development thesis's primary claim is
that the tool is useful to engineers. Three utility properties were
observed during development.

\paragraph{Pipeline automation.} The full path from \Cpp\ source to a
\sysml sequence diagram runs by chaining twelve agent skills
(\cref{tab:skills}), all following a fixed rule set
(\cref{sec:pipeline-overview,sec:evidence-synthesis}). Every skill
produces an auditable intermediate artefact (a CSV, the final
\texttt{.sysml} file, or \texttt{elements.json}/\texttt{sequence.json}
for synthesis) rather than acting invisibly, and works from the live
schema rather than hardcoded labels, so an engineer can inspect, verify
or override any step without rerunning extraction
(\cref{sec:skill-design}).

\paragraph{Scale.} A single Cypher query extracted 2,574 ordered
interactions from CPSCore -- several working days of manual effort --
and identified the nine caller functions across five files to
instrument; \texttt{csp\_matcher} then wrapped those call sites (three
find/replace rules plus a logging-macro patch) and captured the trace,
with no prior knowledge of the codebase required.

\paragraph{Correctness preservation.} Inserting \texttt{csp\_matcher}
probes did not alter CPSCore's observable behaviour: all 1,425 test
assertions passed after instrumentation, and the tool's reversibility
property (\cref{sec:csp-matcher}) restores the source tree to a
byte-identical pre-instrumentation state.

\section{Threats to Validity}

\subsection{Internal Validity}

\begin{description}
  \item[Reference diagram subjectivity.] Reference diagrams reflect a
    single engineer's understanding per scenario; a different reviewer
    might choose different granularity or lifeline selection.

  \item[Instrumentation coverage.] Targets were selected by inspecting
    the static graph, not coverage analysis; a different entry-point
    selection would likely change the trace and H1/H2 result.
\end{description}

\subsection{External Validity}

\begin{description}
  \item[Single controlled system for H1/H2/H3.] CPSCore is a
    moderate-scale, cleanly modular framework. The practitioner session
    offers encouraging but unblinded, four-participant qualitative
    evidence of transfer, not a repeat of the controlled comparisons;
    applying the full hypothesis tests to a second system remains the
    strongest way to strengthen external validity.

  \item[No scenario exercises complementary blind-spot recovery.] All
    four scenarios are found-in-the-wild,
    strengthening confidence that the false-positive-cancellation
    mechanism (\cref{sec:h1-c4c2}). But also means no scenario shows raw union (C4) recovering two
    genuinely complementary blind spots at once. Whether combining ever needs to be additive rather than a precision filter remains untested and is
    left to future work (\cref{sec:future-work}).

  \item[H3's interaction scorer does not verify attribution.] The
    automated check credits a ground-truth item if its source, target
    and method name appear anywhere in a response, without verifying
    they describe the same interaction -- a gap that mattered on S3,
    where two ground-truth interactions share the bare name
    \texttt{stream} with different sources. Figures reported here were
    corrected by manually reading each response
    (\cref{sec:h3-holistic}); a stricter, structurally-aware matcher is
    left to future work.
\end{description}

\subsection{Construct Validity}

\begin{description}
  \item[Edge-level evaluation metric.] The primary H1 metric measures
    precision/recall over unique $\langle\text{sender},\,
    \text{receiver}\rangle$ pairs, abstracting over message count and
    ordering; a complementary LCS-based execution-order metric
    (\cref{sec:h1}) addresses ordering directly. Neither metric weights
    message content or repetition count, so a diagram with the right
    calls in the right order but the wrong volume would still score
    perfectly on both.
\end{description}
